\documentclass[11pt]{article}
\usepackage{amsmath,amssymb,amsthm}
\usepackage{hyperref}

\newtheorem{theorem}{Theorem}
\newtheorem{lemma}{Lemma}

\newcommand{\diam}{\operatorname{diam}}
\newcommand{\dist}{\operatorname{dist}}

\title{A Quadratic WSPD Obstruction in Nagata Dimension Zero}
\author{Run \texttt{fa\_banach\_001}, agent \texttt{agent\_lane\_08}}
\date{2026-07-19}

\begin{document}
\maketitle

\section*{Source Question}

The source paper is Charles Fefferman and Pavel Shvartsman,
\emph{Sharp finiteness principles for Lipschitz selections: long version},
arXiv:1708.00811. In the introduction, after explaining that the paper's
results lead to efficient-computation questions for Lipschitz selection
problems on finite metric spaces, the authors ask whether the results of
Har-Peled and Mendel on well-separated pair decompositions for doubling
metrics can be extended from doubling metrics to metrics of bounded Nagata
dimension.

This packet gives a negative answer to the direct linear-size WSPD extension:
bounded Nagata dimension alone cannot imply an \(O(n)\)-size WSPD for fixed
separation.

\section*{Definitions}

We use the Nagata condition as stated in the source. A metric space \((X,d)\)
satisfies the Nagata condition with constants \(c\in(0,1]\) and \(C\ge 0\) if
for every \(r>0\) there is a cover of \(X\) by subsets of diameter at most
\(r\), at most \(C+1\) of which meet any ball of radius \(cr\).

For a finite metric space \(X\) and a parameter \(s>0\), an \(s\)-well-separated
pair decomposition, abbreviated \(s\)-WSPD, is a family of pairs
\((A_i,B_i)\) of nonempty disjoint subsets of \(X\) such that every unordered
two-point set \(\{x,y\}\subset X\) is covered by at least one pair, meaning
\(x\in A_i,y\in B_i\) or \(y\in A_i,x\in B_i\), and
\[
  \dist(A_i,B_i)\ge s\max\{\diam A_i,\diam B_i\}.
\]
This is the standard metric form of the WSPD separation condition. If one's
convention asks for each two-point set to be covered exactly once, the lower
bound below remains valid.

\section*{Proof Intuition}

Bounded Nagata dimension controls covers scale by scale, but it does not
control packing. The equilateral metric is the extreme example: its Nagata
dimension is \(0\), because below the unit scale singletons work and at or
above the unit scale the whole space works. But a WSPD with separation
\(s>1\) cannot group two equilateral points on either side of a pair: a set
with two points already has the same diameter as its distance to any disjoint
set. Therefore every WSPD pair must be singleton versus singleton, so all
\(\binom n2\) point pairs must be listed separately.

\begin{lemma}
For each \(n\ge2\), the \(n\)-point equilateral metric space \(E_n\), with
\(d(x,y)=1\) for \(x\ne y\), satisfies the Nagata condition with constants
\(c=1\) and \(C=0\).
\end{lemma}

\begin{proof}
Fix \(r>0\). If \(r<1\), cover \(E_n\) by its singleton subsets. Each has
diameter \(0\), and each ball of radius \(r\) contains only its center, hence
meets exactly one member of the cover. If \(r\ge1\), cover \(E_n\) by the
single set \(E_n\), whose diameter is \(1\le r\). Again every ball meets only
one member of the cover. This proves the Nagata condition with \(c=1\) and
\(C=0\), uniformly in \(n\).
\end{proof}

\begin{lemma}
Let \(s>1\). Every \(s\)-well-separated pair \((A,B)\) in \(E_n\) has
\(\#A=\#B=1\).
\end{lemma}

\begin{proof}
If \(A\) contains two distinct points, then \(\diam A=1\). Since \(A\) and
\(B\) are nonempty and disjoint, \(\dist(A,B)=1\). The \(s\)-separation
condition would require \(1\ge s\), a contradiction. Hence \(A\) is a
singleton. The same argument applies to \(B\).
\end{proof}

\begin{theorem}
For every \(s>1\) and every \(n\ge2\), every \(s\)-WSPD of the \(n\)-point
equilateral metric \(E_n\) has at least \(\binom n2\) pairs. Consequently,
there is no bound of the form \(C(s,c,C)\,n\) for WSPDs on all finite metric
spaces satisfying the Nagata condition with constants \(c,C\).
\end{theorem}

\begin{proof}
By the preceding lemma, each \(s\)-well-separated pair in \(E_n\) is of the
form \((\{x\},\{y\})\) with \(x\ne y\). Such a pair covers exactly the unordered
two-point set \(\{x,y\}\). Since a WSPD must cover all \(\binom n2\) unordered
two-point subsets of \(E_n\), it must contain at least \(\binom n2\) pairs.

The Nagata constants of \(E_n\) are \(c=1\), \(C=0\), uniformly in \(n\). Thus
no WSPD theorem depending only on \(s,c,C\) can give a linear-size bound
\(C(s,c,C)n\) for all finite bounded-Nagata-dimension metrics.
\end{proof}

\section*{Scope and Limitations}

The theorem refutes the direct transfer of the usual fixed-separation,
linear-size WSPD conclusion from doubling metrics to bounded Nagata dimension.
It does not rule out weaker decompositions, randomized or approximate
surrogates, or algorithms for Lipschitz selection that exploit bounded Nagata
dimension without producing a standard WSPD of linear size.

\section*{Novelty Check}

The run indexes were searched for arXiv:1708.00811, the source title,
`Lipschitz selections', `finiteness principle', `WSPD', `Har-Peled Mendel',
`doubling metrics', and `Nagata dimension'. No exact prior packet was found.
A bounded web search on 2026-07-19 for the exact WSPD/Nagata phrases found the
source paper and general Nagata references, but no exact equilateral lower
bound recorded as an answer to this source question.

\section*{References}

\begin{thebibliography}{9}

\bibitem{FeffermanShvartsman}
C. Fefferman and P. Shvartsman,
\emph{Sharp finiteness principles for Lipschitz selections: long version},
arXiv:1708.00811.

\bibitem{CallahanKosaraju}
P. B. Callahan and S. R. Kosaraju,
\emph{A decomposition of multidimensional point sets with applications to
\(k\)-nearest-neighbors and \(n\)-body potential fields},
J. ACM 42 (1995), 67--90.

\bibitem{HarPeledMendel}
S. Har-Peled and M. Mendel,
\emph{Fast construction of nets in low-dimensional metrics and their
applications},
SIAM J. Comput. 35 (2006), no. 5, 1148--1184.

\end{thebibliography}

\end{document}

